% ============================================================================
%  preamble.tex -- shared preamble for the math-only unified_theory_math.tex
%  Packages, theorem environments, and the notation macros. Symbols are
%  identical to the source document; only the prose is domain-neutral (pure
%  math / machine learning, no application-domain interpretation).
% ============================================================================

\usepackage[T1]{fontenc}
\usepackage[utf8]{inputenc}
\usepackage{lmodern}
\usepackage[letterpaper,margin=1in]{geometry}
\usepackage{amsmath,amssymb,amsthm,mathtools}
\usepackage{bm}
\usepackage{booktabs}
\usepackage{longtable}
\usepackage{array}
\usepackage{enumitem}
\usepackage{xcolor}
\usepackage{listings}
\usepackage{graphicx}
\usepackage[colorlinks=true,linkcolor=blue!55!black,citecolor=blue!55!black,%
            urlcolor=blue!55!black]{hyperref}
\usepackage{microtype}

% ---- theorem environments --------------------------------------------------
\theoremstyle{plain}
\newtheorem{theorem}{Theorem}[section]
\newtheorem{lemma}[theorem]{Lemma}
\newtheorem{corollary}[theorem]{Corollary}
\newtheorem{proposition}[theorem]{Proposition}
\theoremstyle{definition}
\newtheorem{definition}[theorem]{Definition}
\newtheorem{remark}[theorem]{Remark}

% ---- a light box for open-question / caveat flags --------------------------
\newcommand{\caveat}[1]{\par\noindent\textbf{\textcolor{red!60!black}{Caveat.}}\ #1\par}

% ---- listings style for short inline code / pseudocode ---------------------
\lstset{basicstyle=\ttfamily\small,breaklines=true,columns=fullflexible,
        keepspaces=true,frame=none,xleftmargin=1em}

% ============================================================================
%  NOTATION MACROS  (bind: symbol -> concept, one per concept)
% ============================================================================

% -- generic operators --
\newcommand{\E}{\mathbb{E}}
\newcommand{\Reals}{\mathbb{R}}
\newcommand{\Tr}{\operatorname{Tr}}
\newcommand{\diago}{\operatorname{diag}}
\newcommand{\KLdiv}{\mathrm{KL}}
\newcommand{\GProc}{\mathcal{GP}}
\newcommand{\Poi}{\mathrm{Poisson}}
\newcommand{\Gauss}{\mathcal{N}}
\newcommand{\ELBO}{\mathcal{L}}          % ELBO  (script L; NOT the Cholesky factor)
\newcommand{\transpose}{^{\mathsf{T}}}

% -- model, kernel, structured prior --
\newcommand{\lat}{\lambda}               % latent GP function lambda(x)
\newcommand{\latz}{\lambda_{0}}          % rate offset
\newcommand{\rate}{f}                    % rate f = exp(A lambda + lambda_0)
\newcommand{\lgr}{g}                     % log-rate g = A lambda + lambda_0
\newcommand{\gain}{A}                    % gain (NOT projection a)
\newcommand{\spk}{y}                     % count observation (synonyms r, n)
\newcommand{\Kf}{K}                      % prior kernel function K(x,x')
\newcommand{\Ktil}{\widetilde{K}}        % inducing gram  K(Ztil,Ztil)
\newcommand{\Cmat}{C}                    % structured (grid) prior covariance
\newcommand{\Csm}{C_{\mathrm{smooth}}}   % RBF smoothing factor
\newcommand{\Kmag}{\mathcal{M}}          % kernel magnitude sqrt(v_x v_x')  (NOT M)
\newcommand{\kang}{\theta}               % kernel angle
\newcommand{\Jang}{J}                    % angular term J(theta)  (never a Jacobian)
\newcommand{\vmag}{v}                    % self-magnitude v_x = x^T C x + sigma_0^2
\newcommand{\loc}{\alpha}                % localization weight alpha_i (NOT diffusion alpha_t)
\newcommand{\bwid}{\beta}                % localized-region half-width (prior; NOT diffusion beta_t)
\newcommand{\smoo}{\rho}                 % smoothness SD (prior; NOT correlation)
\newcommand{\sigz}{\sigma_{0}}           % bias std (sigma_0 = exp(raw_sigma_0))
\newcommand{\rfc}{\bm{\xi}_{0}}          % center of the localized region (eps0x, eps0y)
\newcommand{\pixc}{\bm{\xi}}             % grid coordinate
\newcommand{\Amp}{\mathrm{Amp}}          % amplitude: 6th kernel hyperparameter (softplus, default 1.0, optimized)

% -- variational inference & eigenspace --
\newcommand{\ipt}{\widetilde{z}}         % inducing point (NOT subspace coord)
\newcommand{\Ipt}{\widetilde{Z}}         % inducing-point matrix
\newcommand{\latt}{\widetilde{\lambda}}  % inducing values (= u)
\newcommand{\Mind}{M}                    % number of inducing points (NOT magnitude)
\newcommand{\Ntr}{N}                     % number of training inputs
\newcommand{\kvec}{k}                    % cross-covariance vector k(x)=K(Ztil,x)
\newcommand{\pv}{u}                      % projection vector u(x)=Ktil^{-1} k(x)
\newcommand{\vm}{m}                      % variational mean (inducing space)
\newcommand{\vV}{V}                      % variational covariance
\newcommand{\pmean}{\mu}                 % posterior mean mu(x) of the latent
\newcommand{\pvar}{\sigma^{2}}           % posterior variance sigma^2(x)
\newcommand{\Lchol}{L_{\Ktil}}           % Cholesky factor of Ktil (upright L)
\newcommand{\eb}{B}                      % eigenbasis of Ktil
\newcommand{\evals}{\Lambda}             % eigenvalues (diagonal; NOT lambda)
\newcommand{\nb}{n_{b}}                  % kept eigen-dimension
\newcommand{\mb}{m_{b}}                  % eigenspace mean
\newcommand{\Vb}{V_{b}}                  % eigenspace covariance (dense)
\newcommand{\Ktilb}{\Ktil_{b}}           % diagonal inducing gram in eigenspace
\newcommand{\amat}{a}                    % projection a = K Ktil^{-1}
\newcommand{\Kvec}{\mathrm{Kvec}}        % diagonal prior variance vector
\newcommand{\fbar}{\bar{f}}              % expected rate

% -- training-step helpers --
\newcommand{\gE}{g_{\mathrm{E}}}         % E-step gradient helper (NOT log-rate g)
\newcommand{\GE}{G_{\mathrm{E}}}         % E-step Fisher curvature

% -- utility & information --
\newcommand{\Ustd}{U_{\mathrm{std}}}     % standard utility
\newcommand{\Uda}{U_{\mathrm{DA}}}       % distribution-aware utility
\newcommand{\Hmarg}{H_{\mathrm{marg}}}   % marginal count entropy
\newcommand{\Hnoise}{H_{\mathrm{noise}}} % aleatoric Poisson noise entropy
\newcommand{\Hcond}{H_{\mathrm{cond}}}   % conditional entropy over p(x)
\newcommand{\px}{p(x)}                   % input (data) distribution
\newcommand{\mug}{\mu_{g}}               % log-rate mean = A mu + lambda_0
\newcommand{\varg}{\sigma_{g}^{2}}       % log-rate variance = A^2 sigma^2
\newcommand{\corr}{\rho_{\lambda}}       % posterior latent correlation (NOT smoothness)
\newcommand{\ccond}{c}                   % prior conditional covariance
\newcommand{\scal}{\kappa}               % input scaling factor (x* = kappa x_t)
\newcommand{\varr}{\sigma_{r}^{2}}       % residual/aleatoric variance
\newcommand{\subc}{\zeta}                % subspace coordinate
\newcommand{\LW}{W_{0}}                  % Lambert-W principal branch
\newcommand{\gmode}{\bar{g}_{r}}         % Laplace mode
\newcommand{\rmax}{r_{\max}}             % count-truncation bound
\newcommand{\satf}{\sigma_{f}^{2}}       % saturation-kernel amplitude

% -- diffusion --
\newcommand{\dt}{t}                      % diffusion timestep
\newcommand{\xz}{x_{0}}                  % clean input
\newcommand{\xt}{x_{t}}                  % noisy input
\newcommand{\xhat}{\widehat{x}_{0}}      % Tweedie estimate
\newcommand{\bsched}{\beta_{t}}          % noise-schedule variance (NOT the prior beta)
\newcommand{\aret}{\alpha_{t}}           % signal retention (NOT the localization weight alpha)
\newcommand{\abar}{\bar{\alpha}_{t}}     % cumulative signal retention
\newcommand{\epsth}{\epsilon_{\theta}}   % U-Net noise predictor
\newcommand{\score}{s_{\theta}}          % score function
\newcommand{\gw}{w}                      % guidance scale
\newcommand{\ggrad}{g_{t}}               % guidance gradient (utility gradient)
\newcommand{\greg}{\gamma_{\mathrm{reg}}}% regularization weight (input-optimization approach)
